\documentclass[../../main.tex]{subfiles}
\begin{document}

{\bf [解]}
$f(x)=x^3-ax^2+bx-c$とおく．$f(x)=0$が整数解$x=k$を持つ時，
\begin{align}
c=k(k^2-ak+b) \quad\cdots\text{①} \label{eq:1}
\end{align}
となる．

$a,b,c>0$より$k \le 0$の時は$k^2-ak+b>0$だから右辺は非正，左辺は正となって不適である．
従って$k>0$である．よって$k, k^2-ak+b$はいずれも正の整数だから，$c=1,2,\cdots,6$に対して因数になっている．
この組を書き出すと以下の表を得る．
\begin{table}[htb]
  \centering
 \caption{$(c,k,k^2-ak+b)$のあり得る組み合わせ}
\begin{tabular}{c|l}
$c$ & $(k,k^2-ak+b)$ \\
\hline
1 & $(1,1)$ \\
2 & $(2,1)$，$(1,2)$ \\
3 & $(3,1)$，$(1,3)$ \\
4 & $(4,1)$，$(2,2)$，$(1,4)$ \\
5 & $(5,1)$，$(1,5)$ \\
6 & $(6,1)$，$(3,2)$，$(2,3)$，$(1,6)$
\end{tabular}
\end{table}

これら全ての場合に対して$k^2-ak+b$が実際に表の値になるような$(a,b)$の組み合わせを羅列すると以下のようになる．

\begin{enumerate}
 \item[$(1,1)$の時] $a=b \therefore\ (a,b)=(1,1),(2,2),(3,3),(4,4),(5,5),(6,6)$  \\
 \item[$(2,1)$の時] $-2a+b=-3 \therefore\ (a,b)=(2,1),(3,3),(4,5)$ \\
 \item[$(1,2)$の時] $-a+b=1 \therefore\ (a,b)=(1,2),(2,3),(3,4),(4,5),(5,6)$ \\
 \item[$(3,1)$の時] $-3a+b=-8 \therefore\ (a,b)=(3,1),(4,4)$ \\
 \item[$(1,3)$の時] $-a+b=2 \therefore\ (a,b)=(1,3),(2,4),(3,5),(4,6)$ \\
 \item[$(4,1)$の時] $-4a+b=-15 \therefore\ (a,b)=(4,1),(5,5)$ \\
 \item[$(2,2)$の時] $-2a+b=-2 \therefore\ (a,b)=(2,2),(3,4),(4,6)$ \\
 \item[$(1,4)$の時] $-a+b=3 \therefore\ (a,b)=(1,4),(2,5),(3,6)$ \\
 \item[$(5,1)$の時] $-5a+b=-24 \therefore\ (a,b)=(5,1),(6,6)$ \\
 \item[$(1,5)$の時] $-a+b=4 \therefore\ (a,b)=(1,5),(2,6)$ \\
 \item[$(6,1)$の時] $-6a+b=-35 \therefore\ (a,b)=(6,1)$ \\
 \item[$(3,2)$の時] $-3a+b=-7 \therefore\ (a,b)=(3,2),(4,5)$ \\
 \item[$(2,3)$の時] $-2a+b=-1 \therefore\ (a,b)=(1,1),(2,3),(3,5)$ \\
 \item[$(1,6)$の時] $-a+b=5 \therefore\ (a,b)=(1,6)$
\end{enumerate}
である．したがって，重複する$(a,b,c)=(4,5,2)$を1通りとして数えると，あわせて38通りだから求める確率は
\begin{align}
\dfrac{38}{6^3}=\dfrac{19}{108}
\end{align}
である．

\end{document}
